\section{Discusión}\label{sec:discusion}

\subsection{Contribuciones del trabajo}

El documento aporta cuatro resultados que, en conjunto, no existen en la
literatura de bioconversión con \textit{H.~illucens}:

\begin{enumerate}
  \item \textbf{Un modelo unificado alimento--larvas--ambiente.} El modelo
  de \citet{eriksenDynamicModellingFeed2022} describe la larva en un ambiente
  exógeno; aquí la temperatura, la humedad y los gases del contenedor pasan
  a ser estados acoplados (Ec.~\eqref{eq:estado-completo}), con el lecho como nodo hidrotérmico
  intermedio y la forma unificada ventilado--cerrado de la
  Ec.~\eqref{eq:gas-unificado} como instrumento de medición.
  \item \textbf{Un método de medición híbrido con ventana adaptativa.}
  Los cierres de ventilación se dimensionan desde el propio modelo
  (Ecs.~\eqref{eq:delta-min}--\eqref{eq:delta-max},
  Tabla~\ref{tab:delta-adaptativo}): la literatura usa cierres fijos
  \citep{ermolaevGreenhouseGasEmissions2019} o muestreo continuo sin
  acumulación \citep{parodiBioconversionEfficienciesGreenhouse2020a}.
  \item \textbf{Un diseño experimental que resuelve la identificabilidad
  larva--microbiota.} El testigo sin larvas, la línea base temprana y el
  cierre post-cosecha (Sección~\ref{sec:identificabilidad}) separan las dos
  fuentes de CO$_2$ por diseño, no por suposición; la celda de carga añade
  una restricción de consistencia global independiente
  (Ec.~\eqref{eq:celda-carga}).
  \item \textbf{Un gemelo digital con estimación en línea.} La firma de
  terminación CH$_4$/CO$_2$ (Observación~\ref{rem:tf}), el ajuste adaptativo
  de $\Delta$ y la reconstrucción de estados latentes por PINN
  (Sección~\ref{sec:pinn}) convierten el montaje en un sistema de estimación,
  no solo de registro.
\end{enumerate}

\subsection{Posición frente a la literatura}

\begin{table}[htbp]
  \centering
  \caption{Posición del trabajo frente a los estudios de referencia.}
  \label{tab:posicion}
  \begin{tabular}{llll}
    \hline
    Estudio & Medición & Modelo & Estimación en línea \\
    \hline
    Eriksen (2022) & respirometría de larvas & larva, ambiente exógeno & no \\
    Ermolaev (2019) & cámara cerrada + GC puntual & no & no \\
    Parodi (2020) & cámara abierta continua & balance de masa estático & no \\
    Rossi (2024) & 4 cámaras abiertas, 12 h & interpolación empírica & no \\
    Este trabajo & conmutada, $\Delta$ adaptativo, GC & compartimental acoplado & sí (PINN) \\
    \hline
  \end{tabular}
\end{table}

La Tabla~\ref{tab:posicion} resume la posición: cada estudio aporta una
pieza ---el modelo fisiológico, la validación cromatográfica, el balance de
masa, la dinámica temporal--- y el presente trabajo los integra en un solo
marco donde el modelo y la medición se diseñan uno para el otro.

\subsection{Limitaciones}

\begin{description}
  \item[Réplicas secuenciales.] La confusión corrida--tiempo es inherente al
  equipamiento único (Observación~\ref{rem:paralelizacion}); la variabilidad
  entre lotes queda confundida con la variabilidad temporal del laboratorio,
  y solo se mitiga con metadatos, recalibración y el testigo común.
  \item[Ceguera parcial entre cierres.] La delegación del CH$_4$ al GC y del
  O$_2$ fino al modelo (Observación~\ref{rem:o2-supervision}) implica que
  entre cierres esas dos variables no se observan: un evento anaerobio
  breve entre cierres solo lo vería la alarma MOS.
  \item[Mezcla homogénea y gradientes del lecho.] El supuesto A1 se degrada
  en el escenario $N_0 = 700$ (lecho de $\approx 6$~cm), donde la lectura
  puntual del fondo puede subestimar la temperatura de la zona activa
  (Observación sobre lectura puntual, Sección~\ref{sec:camara-lecho}).
  \item[Población homogénea.] El supuesto A2 ignora el ensanchamiento de la
  distribución de pesos (Observación~\ref{rem:homogeneidad}); las réplicas
  acotan su efecto, pero no lo eliminan.
  \item[Una sola condición.] E5 prueba extrapolación en densidad, no en
  sustrato ni en temperatura: la generalización del modelo a otros residuos
  queda por verificar.
\end{description}

\subsection{Una lección de verificación de fuentes}

Durante la elaboración de este documento se detectó que un estudio antes
citado como respaldo metodológico en este campo ---sobre el efecto de la
humedad en las emisiones durante la bioconversión--- fue retractado por su
editorial. La referencia fue sustituida por un estudio equivalente del grupo
de SLU \citep{lindbergProcessEfficiencyGreenhouse2022}. El episodio deja
una regla metodológica que se adopta en lo sucesivo: toda referencia que
sustente una decisión experimental se verifica antes de citarse (estado de
retracción, existencia de los datos reportados), práctica especialmente
necesaria cuando la bibliografía se gestiona en bases sincronizadas que
pueden conservar entradas retractadas.

\subsection{Trabajo futuro}

Cuatro extensiones quedan planteadas sobre el marco construido: (i) cerrar
los balances de nitrógeno y azufre, que convertirían NH$_3$ y H$_2$S de
variables de supervisión en estados del modelo; (ii) reemplazar el supuesto
de población homogénea por una ecuación de estructura poblacional
(Observación~\ref{rem:homogeneidad}); (iii) la paralelización por
multiplexación de gases (Observación~\ref{rem:paralelizacion}), que
convertiría las réplicas secuenciales en simultáneas; y (iv) la dependencia
térmica explícita de los parámetros fisiológicos (factores $Q_{10}$ o
Arrhenius, supuesto B4), necesaria para operar el gemelo fuera de la zona
de confort. La simulación CFD de la línea neumática, anunciada en la
Sección~\ref{sec:diseno}, es el complemento experimental inmediato para
verificar el supuesto A1.
